% Source: https://artofproblemsolving.com/community/c3h2506943_geometry
% Stable identity: aops-topic-2506943-q01
Reflect $E$ in $D$ to get $E'$, so that $D$ is the midpoint of both $BC$ and $EE'$. Then $BECE'$ is a parallelogram (its diagonals bisect each other), and opposite angles of a parallelogram are equal: $\angle BE'C=\angle BEC$. Hence
\[\angle BE'C+\angle BAC=180^\circ.\]
Since $E$ lies on segment $AD$ on the same side of $BC$ as $A$, its reflection $E'$ lies on the other side of $BC$. So $ABE'C$ is a convex quadrilateral with a pair of supplementary opposite angles, hence cyclic.

\begin{asy}
usepackage("amsmath");
size(5cm);
defaultpen(fontsize(10pt));
pair B=(0,0), C=(sqrt(7),0), A=(0.9,2.3);
pair D=(B+C)/2;
real R=abs(A-D); real k=7/4/abs(A-D)/abs(A-D);
pair Ep=D+k*(A-D), E2=D-k*(A-D);
pair O=extension((A+B)/2,(A+B)/2+rotate(90)*(B-A),(B+C)/2,(B+C)/2+(0,1));
draw(circle(O,abs(O-A)),gray);
draw(A--B--C--cycle,linewidth(0.9));
draw(A--E2,blue);
draw(B--Ep--C,red); draw(B--E2--C,red+dashed);
dot(A); dot(B); dot(C); dot(D); dot(Ep); dot(E2);
label("$A$",A,N); label("$B$",B,W); label("$C$",C,E); label("$D$",D,SE); label("$E$",Ep,W); label("$E'$",E2,S);
\end{asy}

The chords $AE'$ and $BC$ of this circle meet at $D$, so by the intersecting chords theorem (power of the point $D$),
\[DA\cdot DE'=DB\cdot DC=\left(\dfrac{\sqrt7}{2}\right)^2=\dfrac74.\]
As $DE'=DE$, we get $DE\cdot DA=\dfrac74$, so $m+n=7+4=\boxed{11}$.
